\SourceLeaf{383}%
tigkeit, als sie uns einen nicht zurückweisbaren Maasstab für die Genauigkeit der Operationen selbst geben, da sie (Rechnungsfehler bei Seite gesetzt) bis auf unmerkliche Kleinigkeiten nichts anderes sein können, als die Aggregate der Fehler, die bei den astronomischen Längenbestimmungen, den Azimuthalbestimmungen, und den Messungen der Winkel im Dreiecksnetze begangen sind. Man kann freilich diese Einflüsse nicht trennen, allein das Dasein des Gesammtfehlers, unabhängig von den Irregularitäten der Erdfigur, ist eine unleugbare Thatsache, wenn auch die Meinung, die man sonst wohl von der \textit{absoluten}, bei allen drei Geschäften erreichten Genauigkeit gehabt hat, merklich herabgestimmt werden muss. Vermuthlich hat jedes seinen Antheil beigetragen. obwohl wir geneigt sind, die grössere Hälfte den gemessenen Dreieckswinkeln zuzuschreiben. Die meisten Operationen, welche Bestandtheile dieser Vergleichungen sind, finden wir zwar in diesem Werke, aber die Winkelmessungen zwischen Mailand und Superga, die in frühern Jahren von französischen Ingenieurs ausgeführt waren, nur in abgekürzter Form, und schon ausgeglichen, so dass man über den Grad ihrer Genauigkeit gar nicht urtheilen kann; inzwischen finden wir in der \textit{Connaissance des tems} 1829 S.\,288, dass Fehler in den Winkelsummen bis zu \(6''8\) dabei vorkommen. Auch das bei Verbindung der Sternwarten von Mailand nach Turin, in Beziehung auf die Übertragung der Orientirung sehr wesentliche Dreieck, Superga, alte und neue Sternwarte von Turin (S.\,254) scheint nicht ganz mit der erforderlichen Genauigkeit gemessen zu sein. Es wäre sehr zu wünschen, dass zur Aufklärung dieses so wichtigen Gegenstandes, wenigstens so weit von den beiden Sternwarten die Rede ist, eine neue geodaetische Verbindung derselben von den dortigen Astronomen ausgeführt werden möchte, was unter Anwendung von zwei Heliotropen und Benutzung des von beiden Sternwarten sichtbaren Platzes S.\,Bernardo di Fenera äusserst leicht sein würde; insofern dort, wie wohl nicht zu zweifeln ist, auch nur einer der übrigen frühern Zwischenpunkte sichtbar ist, würde die ganze Arbeit bloss die Messung von vier Winkeln nöthig machen.

Eine sehr interessante und verdienstliche Arbeit erhalten wir im achten Abschnitt, eine vollständige Wiederholung der von \textsc{Beccaria} 1762\,.\,.\,.\,1764 ausgeführten Breitengradmessung. Bekanntlich liess sich das Resultat dieser Messung mit den in andern Ländern gemessenen Graden gar nicht in Übereinstimmung bringen; die Krümmung des Bogens zwischen den Endpunkten Mondovi und An%
